\section{Beschreibung des Unternehmens}
\label{sec:unternehmen}

Die Aumann~AG mit Sitz in Beelen wurde 1936 gegr\"undet und ist seit
dem 24.~M\"arz 2017 im Prime Standard der Frankfurter Wertpapierb\"orse
notiert; der Emissionspreis lag bei
\EURje{\IPOKurs}.\quelleWeb{Aumann AG, Investor Relations -- Shares}{quelle:shares}{29.07.2026}
Das Unternehmen entwickelt und baut Produktionsanlagen, im Schwerpunkt f\"ur
den elektrifizierten Antriebsstrang. Zum Ende des Gesch\"aftsjahres 2025
besch\"aftigte der Konzern \num{\MitarbeiterZFV} Mitarbeitende nach
\num{\MitarbeiterZFIV} im Vorjahr.\quelleGB{2025}{93}

\subsection{Handelnde Personen in f\"uhrender Position}
\label{subsec:organe}

Der Vorstand besteht aus zwei Mitgliedern, der Aufsichtsrat aus drei
gew\"ahlten Mitgliedern.\quelleGB{2025}{90}\quelleMerken{p3-gb2025-90}

\begin{table}[htbp]
\centering
\caption{Organe der Aumann~AG im Gesch\"aftsjahr 2025.}
\label{tab:organe}
\begin{tabularx}{\linewidth}{@{}llX@{}}
\toprule
\textbf{Organ} & \textbf{Person} & \textbf{Funktion und weitere Mandate}\\
\midrule
Vorstand & Sebastian Roll & Vorstandsvorsitzender (CEO); Supervisor der
  Aumann Technologies (China) Ltd.\ und Mitglied des Board of Directors der
  Aumann Winding and Automation, Inc.\\
Vorstand & Jan-Henrik Pollitt & Finanzvorstand (CFO); Legal Representative
  der Aumann Technologies (China) Ltd.\ und Mitglied des Board of Directors
  der Aumann Winding and Automation, Inc.\\
\midrule
Aufsichtsrat & Gert-Maria Freimuth & Vorsitzender, Vorsitzender des
  Nominierungsausschusses (im Aufsichtsrat seit 21.~November 2016);
  stellvertretender Vorsitzender des Vorstands der MBB~SE\\
Aufsichtsrat & Christoph Weigler & Stellvertretender Vorsitzender,
  Vorsitzender des Pr\"ufungsausschusses (seit 9.~Februar 2017); General
  Manager bei Uber f\"ur Deutschland, \"Osterreich und die Schweiz\\
Aufsichtsrat & Dr.-Ing.\ Saskia Wessel & Mitglied (seit 8.~Juni 2022);
  Leiterin Produkt- und Produktionstechnologie der
  Fraunhofer-Forschungsfertigung Batteriezelle FFB\\
\midrule
Ersatzmitglied & Dr.\ Christof Nesemeier & Ersatzmitglied des Aufsichtsrats
  seit 8.~Juni 2022; Executive Chairman der MBB~SE\\
\bottomrule
\end{tabularx}
\end{table}

\noindent Bemerkenswert ist die personelle Verbindung zur MBB~SE: sowohl der
Aufsichtsratsvorsitzende als auch das Ersatzmitglied bekleiden
F\"uhrungs\-\"amter bei der Mehrheitsaktion\"arin.\quelleErneut{p3-gb2025-90} Mit
Dr.-Ing.\ Saskia Wessel ist zudem die Batteriezellfertigung fachlich im
Aufsichtsrat vertreten.

\subsection{Eigent\"umerstruktur}
\label{subsec:eigentuemer}

Das Grundkapital ist in \num{\AktienGesamt} auf den Inhaber lautende
St\"uckaktien ohne Nennwert
eingeteilt.\quelleWeb{Aumann AG, Investor Relations -- Shares}{quelle:shares}{29.07.2026}\quelleMerken{p4-web-shares}
Gr\"o{\ss}te Aktion\"arin ist die MBB~SE, die zugleich Mutterunternehmen der
Aumann~AG ist.\quelleGB{2025}{93}

\begin{table}[htbp]
\centering
\caption{Aktion\"arsstruktur der Aumann~AG laut Angaben der Gesellschaft.}
\label{tab:eigentuemer}
\begin{tabular}{@{}lr@{}}
\toprule
\textbf{Aktion\"ar} & \textbf{Anteil}\\
\midrule
MBB~SE                          & \Pct{\AnteilMBB}\\
Streubesitz                     & \Pct{\FreefloatAnteil}\\
Nicht n\"aher aufgeschl\"usselt & \Pct{\RestanteilUngeklaert}\\
\bottomrule
\end{tabular}
\end{table}

\noindent \emph{Anmerkung zur Datenlage:} Die ver\"offentlichten Angaben zu
Mehrheitsaktion\"arin und Streubesitz summieren sich nicht auf 100\,\%; die
Differenz von \Pct{\RestanteilUngeklaert} bleibt unerkl\"art. Zudem tragen
die Angaben ungleiche Bezugsdatum: Der Anteil der MBB~SE datiert auf den
31.~Dezember 2025, beim Streubesitz fehlt ein Datum.\quelleErneut{p4-web-shares}

Eine Position in der Gr\"o{\ss}enordnung der Differenz ist belegt. Die
Gesellschaft hielt am 16.~Juli 2026 \num{\EigeneAktienStueck} eigene Aktien
und damit \Pct{\EigeneAktienAnteil} der
Stimmrechte.\quelleWeb{Aumann AG, Ver\"offentlichung nach \S~40 Abs.~1 Satz~2 WpHG vom 16.07.2026}{quelle:wphg}{29.07.2026}
Sie stammen aus dem R\"uckkaufangebot, das am 14.~Juli 2026 zum angehobenen
Preis von \EURje{\RueckkaufPreisZFVI} je Aktie abgeschlossen
wurde.\quelleWeb{Aumann AG, Investor Relations -- Share buyback}{quelle:rueckkauf}{29.07.2026}
Zum Bilanzstichtag 2025 jedoch, auf den sich der Anteil der MBB~SE bezieht,
hielt die Gesellschaft keine eigenen Aktien.\quelleGB{2025}{29} Die Differenz
l\"asst sich der Gr\"o{\ss}enordnung nach damit zuordnen, stichtagsgenau
aufl\"osen l\"asst sie sich nicht. Sie wird hier ausgewiesen und nicht
rechnerisch gegl\"attet.

Die Angaben der Webseite \"andern sich \"uberdies ohne Kennzeichnung: In der
archivierten Fassung vom 14.~September 2025 lauteten dieselben Felder auf
einen Anteil der MBB~SE von \Pct{\AnteilMBBAlt} und einen Streubesitz von
\Pct{\FreefloatAnteilAlt}, bezogen auf die damalige Aktienzahl von
\num{\AktienGesamtMittel}
St\"uck.\quelleWeb{Internet Archive, archivierte Fassung der Seite Shares vom 14.09.2025}{quelle:archiv}{29.07.2026}
Wer beide Fassungen vermischt, erh\"alt eine Differenz, die allein aus der
zwischenzeitlichen Kapitalherabsetzung r\"uhrt.

Die Aktienzahl wurde in zwei Schritten durch Einziehung zur\"uckgekaufter
Aktien verringert: von \num{\AktienGesamtAlt} auf \num{\AktienGesamtMittel}
und weiter auf \num{\AktienGesamt}
St\"uck.\quelleGB{2025}{52} Der R\"uckkauf und die anschlie{\ss}ende
Einziehung sind damit ein wesentlicher Bestandteil der Kapitalallokation, auf
die Abschnitt~\ref{sec:strategie} zur\"uckkommt.

\subsection{Beurteilung durch die Kapitalm\"arkte}
\label{subsec:kapitalmarkt}

Die Aktie wird an den B\"orsenpl\"atzen XETRA, Frankfurt, D\"usseldorf,
Berlin, Stuttgart und M\"unchen gehandelt; Designated Sponsor ist die
ODDO~BHF Corporates \& Markets
AG.\quelleErneut{p4-web-shares}
Die Kursentwicklung der Jahre 2021 bis 2025 ist in
Abschnitt~\ref{subsec:kursverlauf} und Abbildung~\ref{fig:kursverlauf}
dargestellt. Auf Basis des Jahresschlusskurses 2025 von
\EURje{\KursSilvesterZFV} und der aktuellen Aktienzahl ergibt sich eine
Marktkapitalisierung von rund \MioEUR{\MarktkapZFV}.

Zur Einordnung: Der Emissionspreis von \EURje{\IPOKurs} im Jahr 2017 liegt
deutlich \"uber allen Jahresschlusskursen des hier betrachteten Zeitraums.
Die eigentliche Beurteilung, wie der Kapitalmarkt das Unternehmen \"uber den
Kursverlauf hinweg einsch\"atzt, folgt in Abschnitt~\ref{sec:aktienkurs}.

\subsection{Analystenabdeckung}
\label{subsec:analysten}

Die Gesellschaft nennt zwei H\"auser, die die Aktie
begleiten.\quelleWeb{Aumann AG, Investor Relations -- Shares}{quelle:shares}{29.07.2026}\quelleMerken{p5-web-shares}

\begin{table}[htbp]
\centering
\caption{Analystenabdeckung der Aumann-Aktie.}
\label{tab:analysten}
\begin{tabular}{@{}ll@{}}
\toprule
\textbf{Institut} & \textbf{Analyst:in}\\
\midrule
Berenberg            & Yasmin Steilen\\
EQUI.TS              & Thomas J.\ Schiessle\\
\bottomrule
\end{tabular}
\end{table}

\noindent Die Abdeckung hat sich verringert: In der archivierten Fassung der
Seite vom 14.~September 2025 war zus\"atzlich Hauck \& Aufh\"auser mit dem
Analysten Felix Kruse gef\"uhrt.\quelleWeb{Internet Archive, archivierte Fassung der Seite Shares vom 14.09.2025}{quelle:archiv}{29.07.2026} Die Gesellschaft
weist auf diesen Wegfall nicht hin.

\subsection{Botschaften der Analysten}
\label{subsec:analystenbotschaften}

Die Aumann~AG nennt auf ihrer Investor-Relations-Seite die begleitenden
H\"auser, ver\"offentlicht aber weder deren Empfehlungen noch Kursziele oder
Studien.\quelleErneut{p5-web-shares} Bei boerse-express und boerse.de ist
kein Konsens hinterlegt,\quelleWeb{boerse-express, Aumann Aktie kaufen oder
verkaufen? Analyse \& Prognose 2026}{quelle:boerse-express}{02.08.2026}\quelleWeb{boerse.de,
Empfehlungen zur Aumann Aktie}{quelle:boersede-empfehlungen}{02.08.2026}
stockanalysis.com zeigt jedoch f\"ur zwei Analysten die Einstufung Halten
bei einem Kursziel von im Mittel \EURje{\KurszielKonsensDurchschnitt}
(Spanne \EURje{\KurszielKonsensTief} bis
\EURje{\KurszielKonsensHoch}).\quelleWeb{stockanalysis.com,
Analystenkonsens Aumann AG}{quelle:stockanalysis}{02.08.2026}

Die Einzelaktionen best\"atigen dieses Bild: Berenberg (Yasmin Steilen)
stufte im August 2023 von Halten auf Kaufen hoch (Kursziel
\EURje{\KurszielBerenbergVorHochstufung} auf
\EURje{\KurszielBerenbergHochstufung}),\quelleWeb{finanzen.net, Aumann-Aktie im
Plus: Berenberg stuft Aumann auf 'Buy'
hoch}{quelle:berenberg-aug2023}{02.08.2026} hob im November 2023 auf
\EURje{\KurszielBerenbergAnhebung} an\quelleWeb{StreetInsider, Aumann AG PT
Raised to EUR22 at Berenberg}{quelle:berenberg-nov2023}{02.08.2026} und
stufte im Juli 2024 wegen r\"uckl\"aufiger Neuzulassungen zur\"uck auf
Halten (\EURje{\KurszielBerenbergRueckstufung}).\quelleWeb{MarketScreener,
Berenberg senkt Aumann auf 'Hold' - Ziel runter auf 17
Euro}{quelle:berenberg-jul2024}{02.08.2026} EQUI.TS best\"atigte zuletzt im
Juni 2026 ebenfalls Halten.\quelleWeb{aktiencheck.de, Aumann Aktie: Q1/26
schwach, FY 26 ist herausfordernd}{quelle:equits-jun2026}{02.08.2026}

Belegbar ist hingegen der Abgleich zwischen der Prognose der Gesellschaft und
dem tats\"achlichen Verlauf, der in Abschnitt~\ref{sec:strategie} gef\"uhrt
wird. Festzuhalten bleibt, dass die Aktie trotz einer Marktkapitalisierung von
rund \MioEUR{\MarktkapZFV} weiterhin von zwei H\"ausern begleitet wird.
